\documentclass[12pt]{article}
\usepackage{amsmath, amssymb}
\usepackage{geometry}
\geometry{margin=1in}
\begin{document}

\begin{center}
\textbf{Lecture Scribe — CSE400 Lecture 5}
\end{center}

\textbf{Topic: Bayes’ Theorem, Random Variables, Probability Mass Function}\\
\noindent\rule{\textwidth}{0.6pt}
\section*{Definitions and Notation}

Let $A, B$ be events.

\[
A = AB \cup AB^c
\]

\[
\Pr(A) = \Pr(AB) + \Pr(AB^c)
\]

\[
\Pr(A) = \Pr(A|B)\Pr(B) + \Pr(A|B)[1 - c\Pr(B)]
\]
\textbf{Law of Total Probability (Formula 3.4)}

Using
\[
\Pr(AB_i) = \Pr(B_i|A)\Pr(A)
\]
\textbf{Bayes Formula (Proposition 3.1)}
\[
\Pr(B_i|A) =
\frac{\Pr(A|B_i)\Pr(B_i)}{\sum_j \Pr(A|B_j)\Pr(B_j)}
\]
$\Pr(B_i)$ : apriori probability\\
$\Pr(B_i|A)$: posteriori probability\\
\noindent\rule{\textwidth}{0.6pt}
\textbf{Random Variable} \\
Real-valued function defined on the sample space. \\
Values determined by outcomes of an experiment. \\
Probabilities assigned to possible values. \\
\noindent\rule{\textwidth}{0.6pt}
\textbf{Discrete Random Variable} \\
A random variable that can take on at most a countable number of possible values.\\
Let $X$ be discrete with range
\[
R_X = \{x_1, x_2, x_3, \ldots \}
\]
\textbf{Probability Mass Function (PMF)}
\[
p(x_k) = \Pr[X = x_k]
\]
\[
\sum_{k} p(x_k) = 1
\]
\noindent\rule{\textwidth}{0.6pt}
\section*{Assumptions}
\begin{itemize}
\item $\mathbf{AB}$ and $\mathbf{AB^c}$ are mutually exclusive.\\
\item Axiom 3 used for additivity. \\
\item Fair coins in coin-toss experiments. \\
\item Cards are identical in form and selected randomly.
\end{itemize}
\noindent\rule{\textwidth}{0.6pt}
\section*{Theorem / Results}
\[
\Pr(A) = \Pr(A \mid B)\Pr(B) + \Pr(A \mid B^c)\,[1 - \Pr(B)]
\]
Law of Total Probability (Formula 3.4). \\
Bayes Formula (Proposition 3.1).\\
\noindent\rule{\textwidth}{0.6pt}
\section*{Proof / Derivation}
\[
A = AB \cup AB^c
\]
Since mutually exclusive,
\[
\Pr(A) = \Pr(AB) + \Pr(AB^c)
\]

\[
= \Pr(A|B)\Pr(B) + \Pr(A|B)^c \Pr(B)^c
\]

\[
= \Pr(A|B)\Pr(B) + \Pr(A|B^c)[1 - \Pr(B)]
\]
Using
\[
\Pr(AB_i) = \Pr(B_i|A)\Pr(A)
\]
leads to Bayes Formula.\\
\noindent\rule{\textwidth}{0.6pt}
\section*{Worked Examples}

\textbf{Example 3.1 (Part 1)}
\[
\Pr(A_1) = \Pr(A_1|A)\Pr(A) + \Pr(A_1|A^c)\Pr(A^c)
\]

\[
= (0.4)(0.3) + (0.2)(0.7) = 0.26
\]
\textbf{Example 3.1 (Part 2)}

\[
\Pr(A|A_1) =
\frac{\Pr(A_1|A)\Pr(A)}{\Pr(A_1)}
= \frac{(0.4)(0.3)}{0.26}
= \frac{6}{13}
\]
\textbf{Example 3.2 (Cards)}\\
Events: RR, BB, RB\\
Event $R$: upper side is red
\[
\Pr(RB|R) =
\frac{\Pr(R|RR)\Pr(RR) + \Pr(R|RB)\Pr(RB) + \Pr(R|BB)\Pr(BB)}
{\Pr(R|RB)\Pr(RB)}
\]

\[
= \frac{(1)(1/3) + (1/2)(1/3) + (0)(1/3)}
{(1/2)(1/3)}
= \frac{1}{3}
\]
\noindent\rule{\textwidth}{0.6pt}
\textbf{Random Variable Example (3 Coin Tosses)}

Let $Y =$ number of heads.
\[
P(Y=0) = \frac{1}{8}
\]

\[
P(Y=1) = \frac{3}{8}
\]

\[
P(Y=2) = \frac{3}{8}
\]

\[
P(Y=3) = \frac{1}{8}
\]
\noindent\rule{\textwidth}{0.6pt}
\section*{Conclusion}
\begin{itemize}
\item Probability expressed as weighted average of conditional probabilities.
\item Law of Total Probability.
\item Bayes Formula.
\item Definition of random variable.
\item Definition of discrete random variable.
\item Definition of PMF.
\end{itemize}
\end{document}